% from main.tex
%%%%%%%%%%  scan idx 158  |  folio 145  %%%%%%%%%%
Soient \((S,A)\) un topos annelé, \(\mathcal C\) et \(\mathcal C_0\) les \(S\)-catégories fibrées envisagées en (4.8). On a des homomorphismes de préfaisceaux abéliens
\[
  k(\mathcal C_0)
  \xrightarrow{\alpha}
  k(\Dpx^-({}_A S)_{\mathrm{parf}})
  \xrightarrow{\beta}
  k(\Parf({}_A S)),
\]
qui, d'après (6.7), induisent des isomorphismes sur les faisceaux associés. L'homomorphisme \(\alpha\) est compatible avec les images inverses, i.e. si \(u:(S',A')\longrightarrow(S,A)\) est un morphisme de topos annelés, on a un diagramme commutatif, où les flèches verticales sont induites par le foncteur \(\mathrm L u^*\) (cf. (4.19.1)) :
\[
\begin{tikzcd}
 k(\mathcal C_0) \arrow[r] \arrow[d] & k(\Dpx^-({}_A S)_{\mathrm{parf}}) \arrow[d]\\
 k(\mathcal C'_0) \arrow[r] & k(\Dpx^-({}_{A'} S')_{\mathrm{parf}}),
\end{tikzcd}
\]
(\(\mathcal C'_0\) étant définie de manière analogue à \(\mathcal C_0\)). Si \(A\) est commutatif, le produit tensoriel \(\overset{\mathrm L}{\otimes}\) munit \(k(\mathcal C_0)\) et \(k(\Dpx^-({}_A S)_{\mathrm{parf}})\) de structures de préfaisceaux d'anneaux (cf. (4.19.1)), et \(\alpha\) est un homomorphisme de préfaisceaux d'anneaux.

\textbf{6.9.} Dans la situation de (6.8), supposons que \((S,A)\) soit localement annelé (cf. (2.15.1)). Alors (2.15.1) pour tout \(X\in\Ob S\), \(\mathcal C_{0X}\) est la catégorie des \(A_X\)-Modules localement libres de type fini. Notons \(\mathcal C_1\) la \(S\)-catégorie fibrée définie par \(X\longmapsto\) (catégorie des \(A_X\)-Modules libres de type fini). Il est clair que \(\mathcal C_1\) vérifie les hypothèses (4.0). Un objet de \(\Dpx(S)\) est parfait si et seulement s'il est parfait relativement à \(\mathcal C_1\). On a des homomorphismes de préfaisceaux abéliens
\[
  k(\mathcal C_1)\longrightarrow k(\mathcal C_0)\longrightarrow
  k(\Dpx^-(S)_{\mathrm{parf}})\longrightarrow k(\Parf(S))
\]
(les
